\section{Conclusion}
\label{S:conclusion}

We studied an extractive funnel that localizes candidate regions, detects author-stated gap
sentences, and optionally verifies candidates with an LLM. In the reconstructed 40-article SAL
stream, localization retains 203 of 205 labelled limitations while discarding 76.3\% of sentences.
For the saved PubMedBERT predictions, selective verification raises precision from 0.460 to 0.580
and $F_1$ from 0.604 to 0.658, at a recall decrease from 0.878 to 0.761. These results answer RQ1
for the recorded configuration rather than for an independently frozen test protocol.

For RQ2, classical baselines, detector-development studies and auxiliary tasks show that
candidate load, domain, annotation target and reference coverage affect the operating point.
Group-partition localization summaries, alternative silver references and positional sensitivity
scores support narrower conclusions than a single pooled accuracy figure. A literature-gap label policy must not be silently
replaced by an authors'-own limitation policy. For RQ3, accounting on the 40 articles gives a
3.74-fold reduction in estimated LLM input; the selected ten-article baseline study records higher
funnel recall and $F_1$, but does not establish a general advantage over all full-article LLM methods.

The branch additionally supplies a model-labelled reference-construction account and a blind
human-validation interface. Their existence does not establish completed human validation.
The shared prompt-wrapper conflict should be resolved and the revised request frozen before
new inference is evaluated. The next step is an independent evaluation with frozen configurations, complete training and
prediction provenance, adjudicated errors, and measured total processing cost. Reliable source
extraction may support later idea-generation systems, but their novelty and usefulness require
separate experiments.
